% Chapter introduction: why the topic arose, who first asked, what it answers.
% STATUS: written 2026-11-05 (sources of the chapter read in full, each from the page)
Group rings came with a question about derivations from the start. Derivations of
group algebras of torsion-free nilpotent groups were studied by 1978
\cite[p.~434]{FerreroGiambrunoPolcinoMilies1995}. From then on the field asks one
question in many settings: which derivations of a group ring or a group algebra are
inner. Spiegel answered it for the integral group ring of a finite group in 1994
\cite[p.~2955]{Spiegel1994}. Johnson asked it for the Banach algebra $L^{1}(G)$ of
a locally compact group \cite[p.~221]{Losert2008}. Each answer says where a new,
outer derivation can live.

This chapter follows the question through four settings.
Section~\ref{sec:group-hopf:group-algebras} treats the group ring itself. The
integral group ring of a finite group has only inner derivations there. An abelian
group splits off a free part, and the free part measures the outer derivations. An
infinite direct sum of finite $p$-groups has an outer derivation over every field.
The derivation algebra is graded by the quotient of the group by a normal
subgroup with abelian quotient
\cite[Theorem~1, pp.~1, 8]{ArutyunovZhiltsov2023}. Over a finite group the
question becomes a block question, and the first Hochschild cohomology of a
block with nontrivial defect group is often nonzero
\cite[Theorem~A, p.~2]{AntrobusBriggsRubioDegrassi2026}.
Section~\ref{sec:group-hopf:hopf} moves to a Hopf algebra, where a derivation meets
the coproduct. The twisted derivations of Davydov are shown to be symmetries of the
category of representations, taken up to first order. They are computed for a
universal enveloping algebra and for the four-dimensional Hopf algebra.
Section~\ref{sec:group-hopf:quantum} turns to quantum groups, where Woronowicz
builds the differential calculus. A derivation into a bimodule is the basic object,
and the calculi on a compact matrix pseudogroup are classified by right ideals.
The book defines each technical term where it first appears.
Section~\ref{sec:group-hopf:crossed-products} ends with crossed homomorphisms, the
group-map form of a skew derivation, and with the theorem of Losert that every derivation
of $L^{1}(G)$ is inner. The thread through all four sections is the same map with a
new symmetry attached at each step.

\hbfigure{13-group-hopf/laurent}{The group algebra $k[\mathbb{Z}] = k[t,t^{-1}]$. The homomorphism $n\mapsto n$ from $\mathbb{Z}$ to $(k,+)$ gives the derivation $D(t^n)=n\,t^n$, that is $D = t\,d/dt$.}{fig:group:laurent}
